
\documentclass[aspectratio=169,11pt]{beamer}

% =========================================
% PAQUETES
% =========================================
\usepackage[utf8]{inputenc}
\usepackage[T1]{fontenc}
\usepackage[spanish]{babel}
\usepackage{lmodern}
\usepackage{amsmath,amssymb}
\usepackage{tikz}
\usepackage{xcolor}
\usepackage{array}
\usepackage{siunitx}

% TikZ antes de begin document
\usetikzlibrary{
 arrows.meta,
 calc,
 angles,
 quotes,
 positioning,
 decorations.pathreplacing
}

\sisetup{
 output-decimal-marker={,},
 per-mode=symbol
}

% =========================================
% COLORES
% =========================================
\definecolor{azul}{RGB}{0,95,175}
\definecolor{azuloscuro}{RGB}{20,55,100}
\definecolor{verde}{RGB}{0,145,90}
\definecolor{rojo}{RGB}{210,45,65}
\definecolor{naranja}{RGB}{245,100,20}
\definecolor{morado}{RGB}{115,50,170}
\definecolor{gris}{RGB}{245,248,252}

% =========================================
% ESTILO
% =========================================
\usetheme{Madrid}

\setbeamercolor{structure}{fg=azuloscuro}
\setbeamercolor{frametitle}{
 fg=white,bg=azuloscuro
}
\setbeamercolor{block title}{
 fg=white,bg=azul
}
\setbeamercolor{block body}{
 fg=black,bg=blue!5
}
\setbeamercolor{block title example}{
 fg=white,bg=verde
}
\setbeamercolor{block body example}{
 fg=black,bg=green!5
}

\setbeamertemplate{navigation symbols}{}
\setbeamertemplate{footline}[frame number]

% =========================================
% COMANDOS
% =========================================
\newcommand{\resp}[1]{
 \[
 \boxed{\color{morado}{#1}}
 \]
}

\newcommand{\form}[1]{
 \[
 \boxed{#1}
 \]
}

\tikzset{
 vector/.style={
  -{Stealth[length=3mm]},
  very thick
 },
 eje/.style={
  -{Stealth[length=2mm]},
  gray
 }
}

% =========================================
% PORTADA
% =========================================
\title{Magnitudes escalares y vectoriales}
\subtitle{
Representación, operaciones y componentes
}
\author{Clases de Matemática Paraíso}
\date{Física -- Educación Diversificada}

\begin{document}

\begin{frame}
\titlepage
\end{frame}

% =========================================
% OBJETIVOS
% =========================================
\begin{frame}{Objetivos de aprendizaje}

Al finalizar la unidad, el estudiante podrá:

\begin{enumerate}
\item Diferenciar magnitudes escalares
y vectoriales.

\item Representar vectores gráficamente.

\item Identificar magnitud, dirección
y sentido de un vector.

\item Sumar y restar vectores mediante
métodos gráficos.

\item Resolver vectores perpendiculares.

\item Descomponer vectores en componentes
rectangulares.

\item Calcular la magnitud y dirección
de una resultante.
\end{enumerate}

\end{frame}

% =========================================
% SECCIÓN 1
% =========================================
\section{Magnitudes escalares y vectoriales}

\begin{frame}{Magnitudes físicas}

\begin{block}{Concepto}
Una magnitud física es una propiedad
de un cuerpo o fenómeno que puede
medirse y expresarse mediante un
valor numérico y una unidad.
\end{block}

Las magnitudes pueden clasificarse en:

\begin{itemize}
\item Magnitudes escalares.
\item Magnitudes vectoriales.
\end{itemize}

\end{frame}

\begin{frame}{Magnitudes escalares}

\begin{block}{Definición}
Una magnitud escalar queda
determinada mediante un valor
numérico y una unidad.
\end{block}

\begin{center}
\renewcommand{\arraystretch}{1.4}
\begin{tabular}{|l|c|}
\hline
Magnitud & Ejemplo\\
\hline
Masa & \(5\,\mathrm{kg}\)\\
\hline
Tiempo & \(20\,\mathrm{s}\)\\
\hline
Temperatura & \(300\,\mathrm{K}\)\\
\hline
Distancia & \(150\,\mathrm{m}\)\\
\hline
Rapidez & \(12\,\mathrm{m/s}\)\\
\hline
\end{tabular}
\end{center}

\end{frame}

\begin{frame}{Magnitudes vectoriales}

\begin{block}{Definición}
Una magnitud vectorial necesita
magnitud, dirección y sentido
para quedar determinada.
\end{block}

Ejemplos:

\begin{itemize}
\item Desplazamiento.
\item Velocidad.
\item Aceleración.
\item Fuerza.
\end{itemize}

\begin{exampleblock}{Ejemplo}
Una fuerza de \(50\,\mathrm{N}\)
dirigida hacia el este.
\end{exampleblock}

\end{frame}

\begin{frame}{Diferencia entre escalar y vector}

\begin{columns}

\column{0.5\textwidth}

\begin{block}{Escalar}
Un automóvil recorre una distancia
de \(100\,\mathrm{m}\).
\end{block}

\[
d=100\,\mathrm{m}
\]

\column{0.5\textwidth}

\begin{block}{Vector}
Un automóvil se desplaza
\(100\,\mathrm{m}\) hacia el este.
\end{block}

\[
\vec d=100\,\mathrm{m}
\text{ hacia el este}
\]

\end{columns}

\vspace{0.5cm}

\centering
\begin{tikzpicture}
\draw[vector,azul] (0,0)--(5,0);
\node[above] at (2.5,0) {
 \(100\,\mathrm{m}\)
};
\node[right] at (5,0) {Este};
\end{tikzpicture}

\end{frame}

% =========================================
% REPRESENTACIÓN
% =========================================
\section{Representación de vectores}

\begin{frame}{Elementos de un vector}

Un vector se representa mediante
una flecha.

\begin{center}
\begin{tikzpicture}[scale=1.1]

\draw[vector,azul]
(0,0)--(5,2);

\fill (0,0) circle (2pt);

\node[below] at (0,0) {Origen};
\node[above] at (5,2) {Extremo};
\node[above] at (2.5,1) {\(\vec A\)};

\draw[dashed] (0,0)--(5,0);

\draw[naranja,thick]
(1,0) arc[start angle=0,end angle=21.8,
radius=1];

\node[naranja] at (1.5,0.3) {\(\theta\)};

\end{tikzpicture}
\end{center}

\begin{itemize}
\item Magnitud: longitud de la flecha.
\item Dirección: orientación angular.
\item Sentido: hacia dónde apunta.
\end{itemize}

\end{frame}

\begin{frame}{Dirección y sentido}

La dirección puede expresarse
mediante un ángulo respecto
a un eje de referencia.

\begin{center}
\begin{tikzpicture}[scale=0.85]

\draw[eje] (-4,0)--(4,0)
node[right] {\(x\)};

\draw[eje] (0,-2.5)--(0,2.5)
node[above] {\(y\)};

\draw[vector,azul]
(0,0)--(3,1.8);

\draw[vector,verde]
(0,0)--(-3,1.8);

\draw[vector,rojo]
(0,0)--(-3,-1.8);

\draw[vector,naranja]
(0,0)--(3,-1.8);

\node[azul] at (3,2.1) {I};
\node[verde] at (-3,2.1) {II};
\node[rojo] at (-3,-2.1) {III};
\node[naranja] at (3,-2.1) {IV};

\end{tikzpicture}
\end{center}

\end{frame}

\begin{frame}{Vectores iguales y opuestos}

\begin{block}{Vectores iguales}
Tienen la misma magnitud,
dirección y sentido.
\end{block}

\[
\vec A=\vec B
\]

\begin{block}{Vectores opuestos}
Tienen igual magnitud y dirección,
pero sentidos contrarios.
\end{block}

\[
\vec B=-\vec A
\]

\begin{center}
\begin{tikzpicture}
\draw[vector,azul] (0,0)--(3,0);
\draw[vector,rojo] (7,0)--(4,0);

\node[above] at (1.5,0) {\(\vec A\)};
\node[above] at (5.5,0) {\(-\vec A\)};
\end{tikzpicture}
\end{center}

\end{frame}

% =========================================
% SUMA GRÁFICA
% =========================================
\section{Suma gráfica de vectores}

\begin{frame}{Vectores consecutivos}

Dos vectores son consecutivos
cuando el extremo del primero
coincide con el origen del segundo.

\begin{center}
\begin{tikzpicture}[scale=0.95]

\coordinate (O) at (0,0);
\coordinate (A) at (3,0);
\coordinate (B) at (4.3,2);

\draw[vector,azul] (O)--(A);
\draw[vector,verde] (A)--(B);
\draw[vector,naranja] (O)--(B);

\node[below] at (0,0) {Inicio};
\node[above] at (4.3,2) {Final};

\node[below] at (1.5,0) {\(\vec A\)};
\node[right] at (3.8,1) {\(\vec B\)};
\node[above] at (2,1.2) {\(\vec R\)};

\end{tikzpicture}
\end{center}

\form{\vec R=\vec A+\vec B}

\end{frame}

\begin{frame}{Método del triángulo}

\begin{enumerate}
\item Dibujar el primer vector.
\item Colocar el segundo desde
el extremo del primero.
\item Dibujar la resultante desde
el origen inicial hasta el extremo final.
\end{enumerate}

\begin{center}
\begin{tikzpicture}[scale=0.9]

\coordinate (O) at (0,0);
\coordinate (A) at (4,0);
\coordinate (B) at (5,2.5);

\draw[vector,azul] (O)--(A);
\draw[vector,verde] (A)--(B);
\draw[vector,naranja] (O)--(B);

\node[below] at (2,0) {\(\vec A\)};
\node[right] at (4.6,1.3) {\(\vec B\)};
\node[above] at (2.5,1.3) {\(\vec R\)};

\end{tikzpicture}
\end{center}

\end{frame}

\begin{frame}{Vectores concurrentes}

\begin{block}{Definición}
Dos o más vectores son concurrentes
cuando sus líneas de acción
pasan por un mismo punto.
\end{block}

\begin{center}
\begin{tikzpicture}[scale=0.9]

\draw[eje] (-3,0)--(3,0);
\draw[eje] (0,-2)--(0,2.5);

\draw[vector,azul] (0,0)--(2.5,1.5);
\draw[vector,verde] (0,0)--(-2,2);
\draw[vector,rojo] (0,0)--(-2,-1.5);

\node[azul,right] at (2.5,1.5) {\(\vec F_1\)};
\node[verde,left] at (-2,2) {\(\vec F_2\)};
\node[rojo,left] at (-2,-1.5) {\(\vec F_3\)};

\end{tikzpicture}
\end{center}

\end{frame}

\begin{frame}{Método del paralelogramo}

Los vectores se dibujan
desde un mismo origen.

Se completa el paralelogramo
y su diagonal representa
la resultante.

\begin{center}
\begin{tikzpicture}[scale=0.9]

\coordinate (O) at (0,0);
\coordinate (A) at (4,0.5);
\coordinate (B) at (1.3,2.7);
\coordinate (C) at (5.3,3.2);

\draw[vector,azul] (O)--(A);
\draw[vector,verde] (O)--(B);

\draw[dashed,gray,thick] (A)--(C);
\draw[dashed,gray,thick] (B)--(C);

\draw[vector,naranja] (O)--(C);

\node[below] at (2,0.2) {\(\vec A\)};
\node[left] at (0.6,1.5) {\(\vec B\)};
\node[above] at (3,1.8) {\(\vec R\)};

\end{tikzpicture}
\end{center}

\form{\vec R=\vec A+\vec B}

\end{frame}

\begin{frame}{Resta de vectores}

Restar un vector equivale
a sumar su vector opuesto.

\form{\vec R=\vec A-\vec B}

También:

\[
\vec R=\vec A+(-\vec B)
\]

\begin{center}
\begin{tikzpicture}[scale=0.85]

\coordinate (O) at (0,0);
\coordinate (A) at (4,0);
\coordinate (B) at (1.5,2);
\coordinate (N) at (-1.5,-2);
\coordinate (R) at (2.5,-2);

\draw[vector,azul] (O)--(A);
\draw[vector,verde] (O)--(B);
\draw[vector,rojo] (O)--(N);

\draw[vector,rojo,dashed] (A)--(R);
\draw[vector,naranja] (O)--(R);

\node[above] at (2,0) {\(\vec A\)};
\node[left] at (0.7,1.2) {\(\vec B\)};
\node[left] at (-1,-1.1) {\(-\vec B\)};
\node[right] at (2.5,-1.6) {\(\vec R\)};

\end{tikzpicture}
\end{center}

\end{frame}

% =========================================
% VECTORES PERPENDICULARES
% =========================================
\section{Vectores perpendiculares}

\begin{frame}{Vectores perpendiculares}

Dos vectores son perpendiculares
cuando forman un ángulo de
\(90^\circ\).

\begin{center}
\begin{tikzpicture}[scale=0.9]

\coordinate (O) at (0,0);
\coordinate (A) at (4,0);
\coordinate (B) at (4,3);

\draw[vector,azul] (O)--(A);
\draw[vector,verde] (A)--(B);
\draw[vector,naranja] (O)--(B);

\draw (3.7,0)--(3.7,0.3)--(4,0.3);

\node[below] at (2,0) {\(A\)};
\node[right] at (4,1.5) {\(B\)};
\node[above left] at (2,1.5) {\(R\)};

\end{tikzpicture}
\end{center}

\end{frame}

\begin{frame}{Fórmulas para vectores perpendiculares}

\begin{block}{Magnitud de la resultante}

Por el teorema de Pitágoras:

\[
\boxed{R=\sqrt{A^2+B^2}}
\]

\end{block}

\begin{block}{Dirección}

Si \(A>0\) es horizontal
y \(B>0\) es vertical:

\[
\boxed{
\theta=\tan^{-1}\left(\frac{B}{A}\right)
}
\]

\end{block}

\end{frame}

\begin{frame}{Ejemplo resuelto: vectores perpendiculares}

Dos desplazamientos perpendiculares
tienen magnitudes de
\(30\,\mathrm{m}\) y \(40\,\mathrm{m}\).

Determine la magnitud de la resultante.

\textbf{Datos:}

\[
A=30\,\mathrm{m}
\qquad
B=40\,\mathrm{m}
\]

\textbf{Fórmula:}

\[
R=\sqrt{A^2+B^2}
\]

\textbf{Sustitución:}

\[
R=\sqrt{30^2+40^2}
\]

\[
R=\sqrt{900+1600}
\]

\resp{R=50\,\mathrm{m}}

\end{frame}

\begin{frame}{Ejemplo resuelto: dirección}

Utilizando los datos anteriores:

\[
A=30\,\mathrm{m}
\qquad
B=40\,\mathrm{m}
\]

\textbf{Fórmula:}

\[
\theta=\tan^{-1}\left(\frac{B}{A}\right)
\]

\textbf{Sustitución:}

\[
\theta=\tan^{-1}\left(\frac{40}{30}\right)
\]

\resp{\theta\approx53,13^\circ}

La resultante forma un ángulo
de \(53,13^\circ\) sobre
la horizontal positiva.

\end{frame}

% =========================================
% COMPONENTES RECTANGULARES
% =========================================
\section{Componentes rectangulares}

\begin{frame}{Descomposición de un vector}

Un vector puede descomponerse
en dos componentes perpendiculares.

\begin{center}
\begin{tikzpicture}[scale=0.85]

\draw[eje] (-0.3,0)--(5.5,0)
node[right] {\(x\)};

\draw[eje] (0,-0.3)--(0,3.5)
node[above] {\(y\)};

\draw[vector,azuloscuro]
(0,0)--(4,2.8);

\draw[vector,naranja]
(0,0)--(4,0);

\draw[vector,verde]
(4,0)--(4,2.8);

\draw[dashed,gray]
(0,2.8)--(4,2.8);

\draw[naranja,thick]
(0.9,0) arc[start angle=0,end angle=35,
radius=0.9];

\node[naranja] at (1.3,0.35) {\(\theta\)};
\node[below] at (2,0) {\(A_x\)};
\node[right] at (4,1.4) {\(A_y\)};
\node[above] at (2,1.7) {\(\vec A\)};

\end{tikzpicture}
\end{center}

\end{frame}

\begin{frame}{Fórmulas de componentes}

Cuando el ángulo se mide
desde el eje \(x\) positivo:

\begin{block}{Componente horizontal}

\[
\boxed{A_x=A\cos\theta}
\]

\end{block}

\begin{block}{Componente vertical}

\[
\boxed{A_y=A\sin\theta}
\]

\end{block}

Los signos dependen del cuadrante
en el que se encuentra el vector.

\end{frame}

\begin{frame}{Signos de las componentes}

\begin{columns}

\column{0.6\textwidth}

\begin{center}
\begin{tikzpicture}[scale=0.75]

\draw[eje] (-3.5,0)--(3.5,0)
node[right] {\(x\)};

\draw[eje] (0,-3)--(0,3)
node[above] {\(y\)};

\draw[vector,azul] (0,0)--(2.3,2);
\draw[vector,verde] (0,0)--(-2.3,2);
\draw[vector,naranja] (0,0)--(-2.3,-2);
\draw[vector,rojo] (0,0)--(2.3,-2);

\node at (2.6,2.4) {\((+,+)\)};
\node at (-2.6,2.4) {\((-,+)\)};
\node at (-2.6,-2.4) {\((-,-)\)};
\node at (2.6,-2.4) {\((+,-)\)};

\end{tikzpicture}
\end{center}

\column{0.4\textwidth}

\begin{center}
\renewcommand{\arraystretch}{1.5}
\begin{tabular}{|c|c|c|}
\hline
Cuadrante & \(A_x\) & \(A_y\)\\
\hline
I & \(+\) & \(+\)\\
II & \(-\) & \(+\)\\
III & \(-\) & \(-\)\\
IV & \(+\) & \(-\)\\
\hline
\end{tabular}
\end{center}

\end{columns}

\end{frame}

\begin{frame}{Ejemplo: componentes rectangulares}

Una fuerza de \(100\,\mathrm{N}\)
forma un ángulo de \(30^\circ\)
con el eje \(x\) positivo.

Determine sus componentes.

\textbf{Datos:}

\[
F=100\,\mathrm{N}
\qquad
\theta=30^\circ
\]

\textbf{Componente horizontal:}

\[
F_x=F\cos\theta
\]

\[
F_x=100\cos30^\circ
\]

\resp{F_x\approx86,60\,\mathrm{N}}

\end{frame}

\begin{frame}{Ejemplo: componente vertical}

Continuando con el ejemplo:

\[
F=100\,\mathrm{N}
\qquad
\theta=30^\circ
\]

\textbf{Fórmula:}

\[
F_y=F\sin\theta
\]

\textbf{Sustitución:}

\[
F_y=100\sin30^\circ
\]

\resp{F_y=50\,\mathrm{N}}

Por lo tanto:

\[
\vec F=
(86,60\hat i+50\hat j)\,\mathrm{N}
\]

\end{frame}

% =========================================
% SUMA POR COMPONENTES
% =========================================
\section{Suma por componentes rectangulares}

\begin{frame}{Método analítico}

Cuando actúan varios vectores,
podemos calcular la resultante
mediante sus componentes.

\begin{block}{Componentes resultantes}

\[
\boxed{R_x=\sum A_x}
\]

\[
\boxed{R_y=\sum A_y}
\]

\end{block}

\begin{block}{Magnitud}

\[
\boxed{R=\sqrt{R_x^2+R_y^2}}
\]

\end{block}

\end{frame}

\begin{frame}{Dirección de la resultante}

La dirección depende de los
signos de las componentes.

\[
\theta_{\mathrm{ref}}=
\tan^{-1}\left|
\frac{R_y}{R_x}
\right|
\]

Este ángulo es el ángulo
de referencia.

\begin{block}{Ángulo medido desde \(+x\)}

\[
\theta=
\begin{cases}
\theta_{\mathrm{ref}}, & \mathrm{I}\\
180^\circ-\theta_{\mathrm{ref}}, & \mathrm{II}\\
180^\circ+\theta_{\mathrm{ref}}, & \mathrm{III}\\
360^\circ-\theta_{\mathrm{ref}}, & \mathrm{IV}
\end{cases}
\]

\end{block}

Si alguna componente es cero,
la dirección se determina
directamente sobre el eje.

\end{frame}

% =========================================
% EJEMPLO 1
% =========================================
\begin{frame}{Ejemplo: tres fuerzas concurrentes}

Sobre un cuerpo actúan:

\[
F_1=100\,\mathrm{N}
\quad\text{hacia el este}
\]

\[
F_2=80\,\mathrm{N}
\quad\text{hacia el norte}
\]

\[
F_3=60\,\mathrm{N}
\quad\text{hacia el oeste}
\]

Determine la magnitud y dirección
de la fuerza resultante.

\begin{center}
\begin{tikzpicture}[scale=0.7]

\draw[eje] (-3.5,0)--(4.5,0);
\draw[eje] (0,-0.5)--(0,3.5);

\draw[vector,azul] (0,0)--(4,0);
\draw[vector,verde] (0,0)--(0,3);
\draw[vector,rojo] (0,0)--(-2.5,0);

\node[below,azul] at (2,0) {\(F_1\)};
\node[right,verde] at (0,2) {\(F_2\)};
\node[below,rojo] at (-1.5,0) {\(F_3\)};

\end{tikzpicture}
\end{center}

\end{frame}

\begin{frame}{Ejemplo: tabla de componentes}

\begin{center}
\renewcommand{\arraystretch}{1.6}

\begin{tabular}{|c|r|r|}
\hline
Fuerza & \(F_x\) (N) & \(F_y\) (N)\\
\hline
\(F_1\) & 100 & 0\\
\hline
\(F_2\) & 0 & 80\\
\hline
\(F_3\) & -60 & 0\\
\hline
\(\sum\) & 40 & 80\\
\hline
\end{tabular}
\end{center}

\[
R_x=100+0-60=40\,\mathrm{N}
\]

\[
R_y=0+80+0=80\,\mathrm{N}
\]

\end{frame}

\begin{frame}{Ejemplo: magnitud resultante}

\textbf{Fórmula:}

\[
R=\sqrt{R_x^2+R_y^2}
\]

\textbf{Sustitución:}

\[
R=\sqrt{40^2+80^2}
\]

\[
R=\sqrt{1600+6400}
\]

\[
R=\sqrt{8000}
\]

\resp{R\approx89,44\,\mathrm{N}}

\end{frame}

\begin{frame}{Ejemplo: dirección resultante}

\[
\theta=\tan^{-1}\left(\frac{R_y}{R_x}\right)
\]

\[
\theta=\tan^{-1}\left(\frac{80}{40}\right)
\]

\resp{\theta\approx63,43^\circ}

Como ambas componentes son
positivas, la resultante se
encuentra en el primer cuadrante.

\begin{block}{Respuesta}
\[
\boxed{
R\approx89,44\,\mathrm{N}
}
\]

Dirección: \(63,43^\circ\)
sobre el eje \(x\) positivo.
\end{block}

\end{frame}

% =========================================
% EJEMPLO 2
% =========================================
\begin{frame}{Ejemplo: cuatro fuerzas concurrentes}

Determine la magnitud y dirección
de la fuerza resultante.

\begin{center}
\begin{tikzpicture}[scale=0.8]

\draw[eje] (-4,0)--(4.5,0)
node[right] {\(x\)};

\draw[eje] (0,-2.8)--(0,3.3)
node[above] {\(y\)};

\draw[vector,azul] (0,0)--(3.3,0);
\draw[vector,verde] (0,0)--(2.7,2.27);
\draw[vector,morado] (0,0)--(-2.8,1.75);
\draw[vector,rojo] (0,0)--(0,-2.4);

\node[below,azul] at (2.2,0) {\(F_1=100\,N\)};
\node[above,verde] at (2.6,2.3) {\(F_2=200\,N\)};
\node[above,morado] at (-2.8,1.8) {\(F_3=100\,N\)};
\node[right,rojo] at (0,-2.1) {\(F_4=40\,N\)};

\draw[naranja,thick]
(0.65,0) arc[start angle=0,end angle=40,
radius=0.65];

\node[naranja] at (1.1,0.35) {\(40^\circ\)};

\draw[naranja,thick]
(-0.75,0) arc[start angle=180,end angle=148,
radius=0.75];

\node[naranja] at (-1.25,0.3) {\(32^\circ\)};

\end{tikzpicture}
\end{center}

\end{frame}

\begin{frame}{Ejemplo: datos y componentes}

\small

\[
F_1=100\,N,\quad \theta_1=0^\circ
\]

\[
F_2=200\,N,\quad \theta_2=40^\circ
\]

\[
F_3=100\,N,\quad \theta_3=148^\circ
\]

\[
F_4=40\,N,\quad \theta_4=270^\circ
\]

\begin{block}{Fórmulas}

\[
F_x=F\cos\theta
\qquad
F_y=F\sin\theta
\]

\end{block}

\end{frame}

\begin{frame}{Ejemplo: cálculo de componentes}

\small

\textbf{Primera fuerza:}
\[
F_{1x}=100\cos0^\circ=100\,N
\]
\[
F_{1y}=100\sin0^\circ=0\,N
\]
\textbf{Segunda fuerza:}
\[
F_{2x}=200\cos40^\circ
\approx153,21\,N
\]
\[
F_{2y}=200\sin40^\circ
\approx128,56\,N
\]
\textbf{Tercera fuerza:}
\[
F_{3x}=100\cos148^\circ
\approx-84,80\,N
\]
\[
F_{3y}=100\sin148^\circ
\approx52,99\,N
\]
\textbf{Cuarta fuerza:}
\[
F_{4x}=0\,N,\qquad F_{4y}=-40\,N
\]

\end{frame}

\begin{frame}{Ejemplo: tabla de componentes}

\begin{center}
\renewcommand{\arraystretch}{1.5}

\begin{tabular}{|c|r|r|}
\hline
Vector & \(F_x\) (N) & \(F_y\) (N)\\
\hline
\(F_1\) & 100,00 & 0,00\\
\hline
\(F_2\) & 153,21 & 128,56\\
\hline
\(F_3\) & -84,80 & 52,99\\
\hline
\(F_4\) & 0,00 & -40,00\\
\hline
\(\sum\) & 168,41 & 141,55\\
\hline
\end{tabular}
\end{center}

\[
R_x\approx168,41\,N
\]

\[
R_y\approx141,55\,N
\]

\end{frame}

\begin{frame}{Ejemplo: resultante}

Utilizando las componentes
sin redondear en los cálculos:

\[
R=\sqrt{R_x^2+R_y^2}
\]

\[
R\approx
\sqrt{(168,41)^2+(141,55)^2}
\]

\resp{R\approx219,99\,N}

Dirección:

\[
\theta=
\tan^{-1}\left(\frac{R_y}{R_x}\right)
\]

\[
\theta\approx
\tan^{-1}\left(\frac{141,55}{168,41}\right)
\]

\resp{\theta\approx40,05^\circ}

\end{frame}

\begin{frame}{Interpretación de la resultante}

\begin{block}{Respuesta}

La fuerza resultante tiene
una magnitud aproximada de:

\[
\boxed{
\color{morado}{R\approx219,99\,N}
}
\]

Su dirección es:

\[
\boxed{
\color{morado}{\theta\approx40,05^\circ}
}
\]

Se encuentra en el primer cuadrante,
medida desde el eje \(x\) positivo.

\end{block}

\end{frame}

% =========================================
% PRÁCTICAS
% =========================================
\section{Práctica}

\begin{frame}{Práctica 1: clasificación}

Clasifique las siguientes magnitudes
como escalares o vectoriales.

\begin{enumerate}
\item Masa de \(15\,\mathrm{kg}\).
\item Fuerza de \(80\,\mathrm{N}\)
hacia el norte.
\item Tiempo de \(25\,\mathrm{s}\).
\item Velocidad de \(12\,\mathrm{m/s}\)
hacia el este.
\item Distancia de \(200\,\mathrm{m}\).
\item Aceleración de
\(3\,\mathrm{m/s^2}\) hacia abajo.
\end{enumerate}

\end{frame}

\begin{frame}{Práctica 2: vectores perpendiculares}

Dos desplazamientos perpendiculares
tienen magnitudes de
\(12\,\mathrm{m}\) y \(16\,\mathrm{m}\).

Determine:

\begin{enumerate}
\item La magnitud de la resultante.
\item El ángulo respecto al
desplazamiento horizontal.
\item La representación gráfica
de ambos vectores.
\end{enumerate}

\end{frame}

\begin{frame}{Práctica 3: componentes}

Una fuerza de \(150\,\mathrm{N}\)
forma un ángulo de \(35^\circ\)
con el eje horizontal positivo.

Determine:

\begin{enumerate}
\item La componente horizontal.
\item La componente vertical.
\item El cuadrante correspondiente.
\item La representación gráfica
del vector y sus componentes.
\end{enumerate}

\end{frame}

\begin{frame}{Práctica 4: segundo cuadrante}

Una fuerza de \(120\,\mathrm{N}\)
forma un ángulo de \(140^\circ\)
medido desde el eje \(x\) positivo.

Determine:

\begin{enumerate}
\item La componente horizontal.
\item La componente vertical.
\item Los signos de las componentes.
\item La representación gráfica.
\end{enumerate}

\end{frame}

\begin{frame}{Práctica 5: fuerzas concurrentes}

Sobre un cuerpo actúan:

\[
F_1=80\,N
\quad\text{hacia el este}
\]

\[
F_2=60\,N
\quad\text{hacia el norte}
\]

\[
F_3=30\,N
\quad\text{hacia el oeste}
\]

Determine:

\begin{enumerate}
\item Las componentes resultantes.
\item La magnitud de la resultante.
\item La dirección.
\item El cuadrante.
\end{enumerate}

\end{frame}

\begin{frame}{Práctica 6: cuatro fuerzas}

Sobre un punto actúan cuatro fuerzas:

\[
F_1=120\,N,\quad 0^\circ
\]

\[
F_2=90\,N,\quad 60^\circ
\]

\[
F_3=70\,N,\quad 150^\circ
\]

\[
F_4=50\,N,\quad 270^\circ
\]

Determine:

\begin{enumerate}
\item Las componentes de cada fuerza.
\item La suma de componentes.
\item La magnitud de la resultante.
\item Su dirección y cuadrante.
\end{enumerate}

\end{frame}

% =========================================
% FORMULARIO
% =========================================
\section{Resumen}

\begin{frame}{Formulario: vectores}

\small

\begin{columns}

\column{0.5\textwidth}

\textbf{Componentes}

\[
A_x=A\cos\theta
\]

\[
A_y=A\sin\theta
\]

\textbf{Resultantes}

\[
R_x=\sum A_x
\]

\[
R_y=\sum A_y
\]

\column{0.5\textwidth}

\textbf{Magnitud}

\[
R=\sqrt{R_x^2+R_y^2}
\]

\textbf{Ángulo de referencia}

\[
\theta_{\mathrm{ref}}=
\tan^{-1}\left|
\frac{R_y}{R_x}
\right|
\]

\textbf{Resta}

\[
\vec A-\vec B=\vec A+(-\vec B)
\]

\end{columns}

\end{frame}

\begin{frame}{Recomendaciones para resolver problemas}

\begin{enumerate}
\item Leer cuidadosamente el enunciado.
\item Identificar las magnitudes.
\item Dibujar los vectores.
\item Establecer los ejes de referencia.
\item Determinar los ángulos.
\item Calcular las componentes.
\item Sumar las componentes.
\item Calcular magnitud y dirección.
\item Interpretar el resultado físico.
\end{enumerate}

\end{frame}

\begin{frame}

\centering

{\Huge\color{azuloscuro}
Vectores
}

\vspace{0.5cm}

{\Large
Fin de la primera unidad
}

\vspace{0.6cm}

\textbf{Clases de Matemática Paraíso}

\end{frame}

\end{document}
